\begin{enumerate}
\item Une homothétie commute avec tout endomorphisme. Réciproquement, soit $f$ un endomorphisme qui commute avec tout endomorphisme. Si $x \neq 0_E$, soit $F$ un supplémentaire de Vect $(x)$ dans $E$ (C'est pour l'existence de ce supplémentaire qu'on suppose $E$ de dimension finie). Soit $p$ la projection sur $\operatorname{Vect}(x)$ parallèlement à $F$. Alors $f \circ p=p \circ f$. On applique cela en $x$, on obtient $p(f(x))=f(x)$, donc $f(x)$ est lié avec $x$. Et ce, pour tout $x$. Il ne reste plus qu'à appliquer la question précédente.
\item On en déduit que le centre de $\mathcal{M}_n(\K)$ est constitué par les homothéties (en utilisant l'isomorphisme canonique entre $\mathcal{M}_n(\K)$ et $\mathcal{L}\left(\K^n\right)$ ).
\item On retrouve ce résultat en considérant une matrice $M$ du centre et en écrivant
$$
\forall(i, j) \in \llbracket 1, n \rrbracket \quad M E_{i, j}=E_{i, j} M
$$
\end{enumerate}
